\documentclass[11pt]{article}
\usepackage[margin=1in]{geometry}
\usepackage[T1]{fontenc}
\usepackage{lmodern}
\usepackage{hyperref}
\usepackage{longtable}
\usepackage{array}
\usepackage{enumitem}
\usepackage{setspace}
\setstretch{1.08}

\title{MC-RayTracer Environment Lighting: Complete Technical Report}
\author{Codebase Audit (CPU+GPU Paths)}
\date{Generated: 2026-04-27}

\begin{document}
\maketitle
\tableofcontents
\newpage

\section{Scope}
This report documents \textbf{all environment-lighting related code paths} in:
\begin{itemize}[leftmargin=1.5em]
  \item \verb|include/environment.h|
  \item \verb|include/EnvironmentLight.h| + \verb|src/EnvironmentLight.cpp|
  \item \verb|include/ILight.h|
  \item \verb|include/LightFactory.h| + \verb|src/LightFactory.cpp|
  \item \verb|include/LightJsonParser.h| + \verb|src/LightJsonParser.cpp|
  \item \verb|src/EnvironmentLightRegistration.cpp|
  \item \verb|include/scene.h|
  \item \verb|include/texture.h|
  \item \verb|include/query.h| + \verb|include/query.cu|
  \item \verb|include/shader.h|
  \item \verb|src/main.cu|
\end{itemize}

It explains:
\begin{itemize}[leftmargin=1.5em]
  \item where environment light comes from,
  \item what each file and each relevant function does,
  \item what runtime behavior you get today,
  \item what is affected in CPU and GPU modes.
\end{itemize}

\section{Big-Picture Architecture}
\subsection{Two parallel environment-light systems currently coexist}
\textbf{System A (active render path):}
\begin{itemize}[leftmargin=1.5em]
  \item Uses \verb|HDRTextureData| in \verb|texture.h|.
  \item Evaluates background and direct env NEE in \verb|query.h|.
  \item Builds env importance CDF in \verb|main.cu| via \verb|BuildEnvImportance(...)|.
  \item Feeds \verb|EnvImportanceData| into CPU \verb|TraceRayIterative(...)|.
\end{itemize}

\textbf{System B (polymorphic OOP path):}
\begin{itemize}[leftmargin=1.5em]
  \item \verb|ILight| interface, \verb|EnvironmentLight| implementation.
  \item Factory/registry/parsing via \verb|LightFactory| and \verb|LightJsonParser|.
  \item Registered by \verb|EnvironmentLightRegistration.cpp|.
  \item Currently used for runtime parse/validation wiring in \verb|main.cu|, but not as the actual sampler in the integrator.
\end{itemize}

\subsection{Energy contributors in final image}
\begin{enumerate}[leftmargin=1.5em]
  \item \textbf{Direct finite lights}: \verb|ShadeDirect(...)| over \verb|scene.lights| (point + directional).
  \item \textbf{Emissive triangles}: explicit NEE + MIS logic in \verb|TraceRayIterative(...)|.
  \item \textbf{Environment (infinite light)}:
  \begin{itemize}
    \item Miss-path contribution.
    \item Explicit env-light NEE strategy (section 7b in \verb|TraceRayIterative|).
  \end{itemize}
\end{enumerate}

\section{Control Flow: Where Environment Light Comes From}
\subsection{JSON sources}
Environment-related parameters are read from:
\begin{itemize}[leftmargin=1.5em]
  \item top-level \verb|environment| block (\verb|scene.h::parse_environment_block|),
  \item top-level \verb|lights[]| entries with \verb|type == "environment_light"| (\verb|scene.h| and \verb|LightJsonParser.cpp|).
\end{itemize}

\subsection{Precedence in current code}
\begin{itemize}[leftmargin=1.5em]
  \item \verb|environment.hdri_path| sets \verb|scene.sky_hdri_path| first.
  \item In \verb|lights[]|, \verb|environment_light| can override:
  \begin{itemize}
    \item \verb|hdri_path| (if enabled),
    \item \verb|intensity_scale -> scene.environment.map_intensity|,
    \item \verb|tint -> scene.environment.tint|,
    \item \verb|rotation_euler_deg[2] -> scene.environment.map_rotation_deg|.
  \end{itemize}
  \item If \textbf{no directional light exists}, fallback directional sun is synthesized from:
  \verb|environment.sun_dir/sun_tint/sun_intensity|.
\end{itemize}

\section{Header-by-Header Responsibilities}
\subsection{\texttt{include/ILight.h}}
Defines stable ABI for polymorphic lights:
\begin{itemize}[leftmargin=1.5em]
  \item \verb|Sample(normal, u, pdf_out)|: sample incoming direction.
  \item \verb|Pdf(wi)|: solid-angle PDF.
  \item \verb|Eval(wi)|: emitted radiance.
  \item \verb|IsDelta()|: marks singular lights.
\end{itemize}

\subsection{\texttt{include/EnvironmentLight.h}}
Declares environment light implementation:
\begin{itemize}[leftmargin=1.5em]
  \item \verb|EnvironmentLightParams|: enabled flag, HDRI path, intensity/tint, importance resolution, Euler rotation.
  \item \verb|EnvironmentLight| class:
  \begin{itemize}
    \item public MIS API from \verb|ILight|,
    \item private map loading, rotation matrices, and 2D importance distribution.
  \end{itemize}
\end{itemize}

\subsection{\texttt{include/LightFactory.h}}
String-to-constructor registry for \verb|ILight| subclasses:
\begin{itemize}[leftmargin=1.5em]
  \item \verb|Register(type, creator)|
  \item \verb|Create(type, params)|
\end{itemize}

\subsection{\texttt{include/LightJsonParser.h}}
Declares non-throwing parser:
\begin{itemize}[leftmargin=1.5em]
  \item \verb|ParseTopLevelLights(root, out_lights)| and summary counts.
\end{itemize}

\subsection{\texttt{include/environment.h}}
Procedural/latlong environment model and helpers:
\begin{itemize}[leftmargin=1.5em]
  \item defines \verb|Environment| struct (sun, sky, atmosphere, tone-map controls, debug modes),
  \item defines environment radiance evaluation and physical sky math,
  \item defines final image post-process helper \verb|FinalizeImagePixel(...)|.
\end{itemize}

\subsection{\texttt{include/texture.h}}
Defines texture data structures used by environment pipeline:
\begin{itemize}[leftmargin=1.5em]
  \item \verb|HDRTextureData|: float HDR map view + tint + intensity + azimuth rotation.
  \item \verb|EnvImportanceData|: GPU-compatible CDF/PMF view for importance sampling.
  \item \verb|LoadHDRTexture(...)| and \verb|LoadTexture(...)|.
\end{itemize}

\subsection{\texttt{include/scene.h}}
Scene parsing and fallback logic:
\begin{itemize}[leftmargin=1.5em]
  \item parses \verb|environment| block into \verb|Scene.environment|,
  \item parses \verb|lights[]|,
  \item handles \verb|environment_light| special entry,
  \item creates fallback directional sun if no directional light exists and environment enabled.
\end{itemize}

\subsection{\texttt{include/query.h}}
Integrator core:
\begin{itemize}[leftmargin=1.5em]
  \item \verb|sampleHDRI(...)|
  \item env MIS helpers (\verb|env_light_eval/pdf/sample_dir| + CDF sampler)
  \item \verb|TraceRayIterative(...)| miss-path env contribution and env NEE strategy.
\end{itemize}

\subsection{\texttt{include/query.cu}}
Render wrapper:
\begin{itemize}[leftmargin=1.5em]
  \item GPU kernel and CPU loop call \verb|TraceRayIterative(...)|,
  \item passes \verb|EnvImportanceData*| only through CPU path today.
\end{itemize}

\subsection{\texttt{include/shader.h}}
Finite direct light shading path:
\begin{itemize}[leftmargin=1.5em]
  \item \verb|ShadeDirect(...)| for point/directional \verb|Light| array,
  \item directional sunlight shadow behavior in \verb|IsInShadow(...)|.
\end{itemize}

\section{Function-by-Function Catalog}
\subsection{\texttt{include/environment.h}: full environment utility and sky model}
\textbf{Enum/string helpers}
\begin{itemize}[leftmargin=1.5em]
  \item \verb|EnvironmentPresetToString(...)|: preset enum to token.
  \item \verb|EnvironmentPresetFromString(...)|: token to preset enum.
\end{itemize}

\textbf{Math primitives}
\begin{itemize}[leftmargin=1.5em]
  \item \verb|env_clampf|: clamp scalar.
  \item \verb|env_lerp_f|: scalar lerp.
  \item \verb|env_lerp|: vector lerp.
  \item \verb|env_smoothstep|: smooth interpolation.
  \item \verb|env_fract|: fractional part.
  \item \verb|env_safe_normalize|: robust normalization with fallback.
  \item \verb|env_hash3|: deterministic noise hash (stars/twinkle support).
\end{itemize}

\textbf{Direction-angle conversion}
\begin{itemize}[leftmargin=1.5em]
  \item \verb|DirectionFromElevationAzimuthDeg|: elevation/azimuth to world direction.
  \item \verb|ElevationAzimuthFromDirectionDeg|: inverse mapping.
\end{itemize}

\textbf{Preset application}
\begin{itemize}[leftmargin=1.5em]
  \item \verb|ApplyEnvironmentPreset|: fills environment defaults by preset; custom exits early preserving JSON.
\end{itemize}

\textbf{Lat-long map sampling}
\begin{itemize}[leftmargin=1.5em]
  \item \verb|env_sample_latlong_texel|: fetches nearest texel from 8-bit map.
  \item \verb|EvaluateLatLongEnvironment|: bilinear sampled latlong evaluation with yaw rotation and intensity.
\end{itemize}

\textbf{Atmosphere/physics helpers}
\begin{itemize}[leftmargin=1.5em]
  \item \verb|env_earth_radius_m|, \verb|env_beta_R_sea_level|, \verb|env_beta_M_sea_level|,
  \item \verb|env_turbidity_to_mie_factor|,
  \item \verb|env_ray_sphere|,
  \item \verb|env_phase_rayleigh|, \verb|env_phase_mie_hg|.
\end{itemize}

\textbf{Scattering and celestial terms}
\begin{itemize}[leftmargin=1.5em]
  \item \verb|env_view_transmittance|: Beer-Lambert transmittance through atmosphere.
  \item \verb|env_single_scattering|: Nishita-style single scattering (Rayleigh + Mie).
  \item \verb|env_sun_disk|: physically motivated bright solar disc term.
  \item \verb|env_legacy_sky|: artistic non-physical sky fallback.
  \item \verb|env_night_sky|: moon + stars model.
\end{itemize}

\textbf{Output-domain helpers}
\begin{itemize}[leftmargin=1.5em]
  \item \verb|env_aces_channel| and \verb|env_tonemap_aces|: display tonemapping.
  \item \verb|EvaluateEnvironment|: master linear radiance environment function.
  \item \verb|FinalizeImagePixel|: exposure + optional ACES applied once at end.
  \item \verb|EnvironmentDebugPrint|: debug dump of resolved parameters.
\end{itemize}

\subsection{\texttt{src/EnvironmentLight.cpp}: OOP environment light implementation}
\textbf{Internal helpers (anonymous namespace)}
\begin{itemize}[leftmargin=1.5em]
  \item \verb|clamp01|: safe [0,1] clamp.
  \item \verb|luminance709|: BT.709 luminance.
  \item \verb|mat_mul|, \verb|mat_transpose|, \verb|mat_mul_vec|: 3x3 rotation math.
  \item \verb|dir_from_uv|, \verb|uv_from_dir|: latlong direction/UV conversion.
  \item \verb|sample_cdf|: CDF inversion for importance sampling.
\end{itemize}

\textbf{Class methods}
\begin{itemize}[leftmargin=1.5em]
  \item \verb|EnvironmentLight::EnvironmentLight|: copies params, sets env fields, builds rotation + map + PMF/CDF.
  \item \verb|HasMap|: verifies map availability.
  \item \verb|LoadLatLongMap|: loads map via \verb|stbi_load| (8-bit RGB), stores as \verb|TextureData|.
  \item \verb|BuildRotationMatrices|: Euler XYZ to world/env transform matrices.
  \item \verb|EnvToWorld| / \verb|WorldToEnv|: apply transform.
  \item \verb|EvalLatLong|: evaluate map-backed environment radiance through \verb|EvaluateEnvironment|.
  \item \verb|BuildImportanceDistribution|: computes row/column CDF and PMF with \verb|sin(theta)| weighting.
  \item \verb|Sample|: sampled direction from PMF/CDF or uniform-sphere fallback.
  \item \verb|Pdf|: solid-angle PDF from PMF and pixel differential solid-angle.
  \item \verb|Eval|: emitted radiance for world direction.
  \item \verb|IsDelta|: returns false.
\end{itemize}

\subsection{\texttt{include/LightFactory.h} + \texttt{src/LightFactory.cpp}}
\begin{itemize}[leftmargin=1.5em]
  \item \verb|Register(type, creator)|: thread-safe insertion/replacement in registry map.
  \item \verb|Create(type, params)|: thread-safe lookup and instance creation.
\end{itemize}

\subsection{\texttt{include/LightJsonParser.h} + \texttt{src/LightJsonParser.cpp}}
\textbf{Utility readers}
\begin{itemize}[leftmargin=1.5em]
  \item \verb|read_number|, \verb|read_bool|, \verb|read_string|, \verb|read_vec3|, \verb|read_res2|.
\end{itemize}
\textbf{Main parsing}
\begin{itemize}[leftmargin=1.5em]
  \item \verb|parse_environment_light_params|: validates and fills \verb|EnvironmentLightParams|.
  \item \verb|ParseTopLevelLights|: walks \verb|lights[]|, creates \verb|environment_light| via factory, returns loaded/skipped/failed counters.
\end{itemize}

\subsection{\texttt{src/EnvironmentLightRegistration.cpp}}
\begin{itemize}[leftmargin=1.5em]
  \item static \verb|kRegisteredEnvironmentLight|: registers type key \verb|"environment_light"| with factory lambda at startup.
\end{itemize}

\subsection{\texttt{include/texture.h} environment-related functions}
\begin{itemize}[leftmargin=1.5em]
  \item \verb|HDRTextureData::sample(u,v)|: bilinear float HDR sample.
  \item \verb|LoadHDRTexture(path)|: \verb|stbi_loadf| loader for linear HDR.
  \item \verb|EnvImportanceData|: POD view for PMF/CDF used by env sampling.
\end{itemize}

\subsection{\texttt{include/scene.h} environment-related functions}
\begin{itemize}[leftmargin=1.5em]
  \item \verb|parse_environment_block|: loads all \verb|environment| keys.
  \item \verb|parse_one_light|: parses point/area/directional finite lights.
  \item \verb|lights[] environment_light branch|: maps env-light JSON into \verb|scene.environment| and \verb|scene.sky_hdri_path|.
  \item fallback sun block: if no directional light and \verb|environment.enabled|, synthesize directional sun.
  \item \verb|LoadSceneFromFile|: parses JSON text and dispatches to \verb|parse_scene|.
\end{itemize}

\subsection{\texttt{include/query.h} environment-related functions}
\begin{itemize}[leftmargin=1.5em]
  \item \verb|sampleHDRI|: world direction to equirectangular sample.
  \item \verb|env_light_eval|: rotated env evaluation with \verb|tint * intensity|.
  \item \verb|env_cdf_sample|: device-safe binary CDF inversion.
  \item \verb|env_light_pdf|: PMF-to-solid-angle PDF conversion, fallback \verb|1/(4pi)|.
  \item \verb|env_light_sample_dir|: 2-stage CDF inversion (row then column), fallback uniform sphere.
  \item \verb|TraceRayIterative(..., hdri, env_importance)|:
  \begin{itemize}
    \item miss path: adds environment radiance, MIS-weights BSDF strategy in \verb|nee_mode==2|,
    \item section 7b: explicit env-light NEE sampling with power heuristic against BRDF pdf.
  \end{itemize}
\end{itemize}

\subsection{\texttt{include/query.cu} environment-related functions}
\begin{itemize}[leftmargin=1.5em]
  \item \verb|renderBatchCUDA|: GPU kernel calls \verb|TraceRayIterative| (without env importance pointer today).
  \item \verb|render(...)| host wrapper:
  \begin{itemize}
    \item CPU branch passes \verb|env_importance| to \verb|TraceRayIterative|,
    \item GPU branch does not pass env-importance data into kernel in current implementation.
  \end{itemize}
\end{itemize}

\subsection{\texttt{include/shader.h} environment-light interaction}
\begin{itemize}[leftmargin=1.5em]
  \item \verb|IsInShadow|: directional-light shadow rays (infinite distance test).
  \item \verb|ShadeDirect|: accumulates finite \verb|scene.lights| direct contribution.
\end{itemize}

\subsection{\texttt{src/main.cu} environment pipeline functions}
\begin{itemize}[leftmargin=1.5em]
  \item \verb|BuildEnvImportance(const HDRTextureData&, float az_rot)|:
  \begin{itemize}
    \item computes luminance-weighted PMF and CDFs over HDR map,
    \item uses \verb|sin(theta)| Jacobian in sampleHDRI UV convention,
    \item writes \verb|EnvImportanceData| view.
  \end{itemize}
  \item scene load path stores \verb|scene_json_path| then calls:
  \begin{itemize}
    \item \verb|ParseTopLevelLights| (runtime parse/validation and factory wiring),
    \item \verb|BuildEnvImportance| for CPU render.
  \end{itemize}
  \item HDRI setup:
  \begin{itemize}
    \item sets \verb|tint|, \verb|intensity|, \verb|az_rot| into HDR data structs.
  \end{itemize}
\end{itemize}

\section{How It Works End-to-End at Runtime}
\subsection{Step 1: Parse}
\begin{enumerate}[leftmargin=1.5em]
  \item JSON parsed by \verb|SceneIO::LoadSceneFromFile|.
  \item \verb|environment| block fills \verb|Scene.environment|.
  \item \verb|lights[]| parsed:
  \begin{itemize}
    \item directional/point/area become finite \verb|Light| entries,
    \item \verb|environment_light| updates HDRI path/tint/intensity/rotation controls.
  \end{itemize}
  \item If no directional exists and environment enabled, fallback sun light is added.
\end{enumerate}

\subsection{Step 2: Prepare data in \texttt{main.cu}}
\begin{enumerate}[leftmargin=1.5em]
  \item Load float HDRI via \verb|LoadHDRTexture|.
  \item Copy to:
  \begin{itemize}
    \item GPU struct (\verb|d_hdri|): includes tint/intensity/az\_rot,
    \item CPU struct (\verb|h_hdri_data|): includes tint/intensity/az\_rot.
  \end{itemize}
  \item Build CPU importance map CDF via \verb|BuildEnvImportance|.
  \item Pass \verb|EnvImportanceData*| to CPU render path.
\end{enumerate}

\subsection{Step 3: Integrate in \texttt{TraceRayIterative}}
\textbf{When a ray misses:}
\begin{itemize}[leftmargin=1.5em]
  \item evaluate env radiance \verb|Le = env_light_eval(hdri, dir)|,
  \item in MIS mode for bounced non-delta surface events:
  \begin{itemize}
    \item compute \verb|pdf_env = env_light_pdf(...)|,
    \item apply BSDF-strategy weight \verb|w_bsdf = power_heuristic(prev_pdf, pdf_env)|,
    \item accumulate \verb|throughput * Le * w_bsdf|.
  \end{itemize}
  \item otherwise accumulate full \verb|throughput * Le|.
\end{itemize}

\textbf{At each non-miss surface hit (env NEE strategy 7b, MIS mode):}
\begin{itemize}[leftmargin=1.5em]
  \item sample \verb|wi| from env CDF (\verb|env_light_sample_dir|),
  \item shadow-test to infinity,
  \item evaluate \verb|Le| and BRDF \verb|f|,
  \item get BRDF pdf \verb|pdf_B|,
  \item MIS weight \verb|w_L = power_heuristic(pdf_L, pdf_B)|,
  \item accumulate \verb|throughput * Le * f * (NdotL/pdf_L) * w_L|.
\end{itemize}

\section{What It Affects}
\begin{itemize}[leftmargin=1.5em]
  \item \textbf{Background color}: miss rays now use rotated/tinted/intensity-scaled HDRI.
  \item \textbf{Direct illumination quality}: explicit env NEE reduces variance vs pure BSDF-only escape.
  \item \textbf{Sun behavior}: if no explicit directional, fallback sun follows \verb|environment.sun_*|.
  \item \textbf{Scene tone/warmth}: \verb|environment_light.tint| and \verb|intensity_scale| directly modulate env radiance.
\end{itemize}

\section{Current Limitations / Notes}
\begin{itemize}[leftmargin=1.5em]
  \item GPU kernel path currently does not consume uploaded env importance CDF pointers in \verb|renderBatchCUDA| path; CPU path does.
  \item \verb|environment_light.enabled=false| only disables the light-entry override branch; if \verb|environment.hdri_path| is set, HDRI may still be active.
  \item \verb|rotation_euler_deg| is effectively yaw-only in scene parser (\verb|[2]| used as azimuth), while OOP \verb|EnvironmentLight| can build full Euler matrices.
  \item OOP \verb|EnvironmentLight| is runtime-parsed and registered, but integrator uses the float-HDR/CDF path for actual sampling.
  \item In \verb|query.cu| AOV miss assignment still uses raw \verb|sampleHDRI| in one path (without env helper modifiers).
\end{itemize}

\section{Practical Mental Model}
\begin{enumerate}[leftmargin=1.5em]
  \item \textbf{Scene parser decides controls}: sun params, HDRI path, env tint/intensity/rotation.
  \item \textbf{main.cu materializes resources}: loads HDR, builds PMF/CDF, prepares render inputs.
  \item \textbf{query.h performs physics estimator}: combines BSDF strategy and env-light strategy with MIS.
  \item \textbf{shader.h handles finite lights}: directional sun fallback and point lights remain in direct-light loop.
\end{enumerate}

\section{Conclusion}
You currently have a \textbf{hybrid environment-lighting system}:
\begin{itemize}[leftmargin=1.5em]
  \item robust, physically-motivated environment model (\verb|environment.h|),
  \item integrated CPU env MIS path with rotation/tint/intensity and CDF importance sampling,
  \item finite-light path for sun/point lights,
  \item OOP light architecture compiled and runtime-validated but not yet the active integrator sampling backend.
\end{itemize}

If you want a single unified architecture, the next step is to choose one of:
\begin{itemize}[leftmargin=1.5em]
  \item make OOP \verb|EnvironmentLight| the sole sampler/evaluator in render path, or
  \item keep current float-HDR path as canonical and use OOP path only for tooling/authoring.
\end{itemize}

\end{document}
